% Revised one-page proposal (course format: goal, method, prior work and contribution, resources, challenges and
% feasibility).
\documentclass[10pt]{article}
\usepackage[a4paper,margin=1.5cm]{geometry}
\usepackage{amsmath,xcolor,multicol}
\usepackage[hidelinks]{hyperref}
\usepackage[numbers,sort&compress]{natbib}
\usepackage{enumitem}
\setlist{nosep,leftmargin=1.1em,topsep=1pt}
\setlength{\parindent}{0pt}\setlength{\parskip}{3pt}
\setlength{\bibsep}{0pt}
\pagestyle{empty}
\newcommand{\head}[1]{\textbf{\color[HTML]{1f4e8c}#1.}~}

\begin{document}
\begin{center}
{\large\bfseries Is Optical Flow Needed for Stable Sketch Flipbooks?}\\[2pt]
Final project proposal \textbullet{} Modern Computer Vision\\
Rona Lavi, Ron Libman, Yamit Ehrlich, Sophie Danilov \textbullet{} \url{https://github.com/LaviRona/FlipBook}
\end{center}

\head{Goal} We want to turn a short video into a flipbook of abstract line drawings, 16 strokes per page drawn with
CLIPasso~\citep{vinker2022clipasso}, that is recognizable on every page and stable when flipped. Drawing each frame
independently gives good pages but flickers, because every drawing starts from a new initialization. We ask whether drawing
each frame from the previous drawing moved by optical flow removes the flicker, whether errors accumulate, and which
part of the procedure is responsible.

\head{Method} We compare variants that each add one component to the previous one:
\begin{itemize}
\item \textbf{Per-frame}: CLIPasso from scratch on every frame, 2{,}000 iterations.
\item \textbf{Warm start}: frame $t{+}1$ starts from the strokes of frame $t$ and is refined for 150 iterations.
\item \textbf{Propagation without flow}: the warm start, plus re-initialization of strokes in occluded regions
  (forward--backward check~\citep{sundaram2010dense}), a regularizer $\lambda\lVert\theta-\theta_0\rVert^2$ that keeps
  the strokes near their starting position, and a longer keyframe refinement every $K$ frames.
\item \textbf{Flow propagation (proposed)}: as above, but every control point is first moved by RAFT~\citep{teed2020raft}
  optical flow. The two share one implementation, so at equal $\lambda$ their difference is the effect of the flow.
\end{itemize}
\emph{Measures.} Flicker: the chamfer distance (px) between drawing $t$, warped to frame $t{+}1$, and drawing $t{+}1$.
The warp uses GMFlow~\citep{xu2022gmflow,xu2023unifying}, a flow model that no method uses, so the metric cannot favour flow
propagation. Fidelity: LPIPS~\citep{zhang2018lpips} and CLIP ViT-L/14 similarity to the frame. Cost: iterations per
frame.
\emph{Protocol.} DAVIS~\citep{pont2017davis} clips of 24 frames. We choose $\lambda$ (five values) and
$K\in\{6,12,\infty\}$ on one development clip by a rule fixed in advance (lowest flicker with LPIPS $\le1.10\times$
and CLIP $\ge0.95\times$ per-frame), keep them fixed, and then test on three unseen clips with
three seeds each, reporting 95\% bootstrap intervals.
\emph{Success criteria.} (C1) Flow propagation flickers at least 10\% less than propagation without flow. (C2) Without
keyframes, fidelity decreases over the sequence, and keyframes prevent this. (C3) Flow
propagation reaches flicker $\le0.70\times$ per-frame at LPIPS $\le1.10\times$ and CLIP $\ge0.95\times$, for at most
half the cost.

\head{Prior work and our contribution} CLIPasso optimizes the control points of B\'ezier strokes through the
differentiable rasterizer DiffVG~\citep{li2020diffvg}. Video sketching methods optimize the whole video jointly:
Zheng et al.~\citep{zheng2024sketchvideo} against a neural atlas, Zhu et al.~\citep{zhu2025vector,zhu2026hypersketch}
with smooth control-point trajectories initialized from point tracks. Zheng et al.\ set optical flow aside as ``the
most naive approach'', warning of ``cumulative errors'', without testing it. Moving strokes along the flow frame by
frame is older: painterly rendering did it and was found to flicker as motion errors
accumulate~\citep{litwinowicz1997processing,collomosse2004mid}, while in video style transfer a flow-warped start clearly
helps~\citep{ruder2016artistic}. Our contribution is a controlled test of these starting points for CLIP-optimized line
sketches, including the no-flow control that a direct comparison with per-frame drawing would omit, on unseen clips
with a flicker metric that uses an independent flow model.

\head{Resources} One to four GPUs on the lab's Slurm cluster (RTX PRO 6000, 96\,GB). The full set of about 100 runs of 24 frames over four clips takes about 15
GPU-hours. Nothing is trained: we use pretrained CLIP, U2Net, RAFT and GMFlow and public DAVIS clips.

\head{Challenges and feasibility}
\begin{itemize}
\item \emph{Optical flow is wrong at occlusions and under fast motion} and can tear strokes apart. Re-initialization
  and keyframes address this, and the test clips range from slow motion to a fast, partly occluded jump.
\item \emph{CLIPasso is not deterministic.} We use three seeds, compare methods on the same clip and seed, and report
  bootstrap intervals.
\item \emph{The metric may measure the wrong thing.} Pixel warping error also depends on the amount of ink, so we
  check the metrics against simple controls on the development clip.
\item \emph{Feasibility.} Nothing is trained, and a run takes 3--40 GPU-minutes. If time runs short, we drop the $K$ sweep and then one seed.
\end{itemize}


\vspace{-8pt}
\begin{multicols}{2}
\scriptsize
\bibliographystyle{abbrvnat}
\bibliography{refs}
\end{multicols}
\end{document}
